\documentclass[10pt]{article}
\usepackage{amsmath, amssymb}
\usepackage{geometry}
\usepackage{hyperref}
\usepackage{physics}
\usepackage{cite}

\geometry{margin=1in}

\title{Noise Models for Superconducting Qubits}
\author{Qubit Noise Spectroscopy Project \\ University of Waterloo}
\date{}

\begin{document}

\maketitle

\begin{abstract}
We describe the three canonical noise contributions to the power spectral density $S(\omega)$ of a superconducting qubit: $1/f^\alpha$ noise arising from a distribution of two-level fluctuators, a white noise floor from Johnson-Nyquist thermal fluctuations, and Lorentzian peaks from individual two-level system (TLS) defects. Each model is derived from its physical origin, expressed mathematically, and discussed in the context of the filter function formalism developed in documents 02 and 03.
\end{abstract}

\section{Overview}

The noise power spectral density $S(\omega)$ of a superconducting qubit encodes the full frequency-resolved noise environment experienced by the qubit. In practice, $S(\omega)$ for a superconducting qubit is well described by the sum of three physically distinct contributions:

\begin{equation}
    \boxed{S(\omega) = \frac{A}{\omega^\alpha} + S_0 + \sum_k \frac{A_k \gamma_k^2}{(\omega - \omega_k)^2 + \gamma_k^2}}
\end{equation}

where the first term is $1/f^\alpha$ noise, the second is a white noise floor, and the sum represents contributions from individual TLS defects. Each term dominates in a different frequency regime: $1/f^\alpha$ at low frequencies, white noise at high frequencies, and TLS peaks at specific intermediate frequencies.

\section{\texorpdfstring{$1/f^\alpha$}{1/f} Noise}

\subsection{Physical Origin}

The qubit frequency $\omega_0$ is sensitive to environmental parameters — charge on the island, magnetic flux through the loop, critical current of the Josephson junction. Each of these is perturbed by fluctuators: charge traps, magnetic impurities, and atomic tunneling systems in the substrate and oxide layers. Each fluctuator switches randomly between two states at a characteristic rate $\gamma_k$, producing a Lorentzian noise contribution:

\begin{equation}
    S_k(\omega) \propto \frac{\gamma_k}{\omega^2 + \gamma_k^2}
\end{equation}

When many fluctuators are present with a broad distribution of switching rates, the total noise is obtained by summing over all fluctuators. If the switching rates are distributed uniformly on a logarithmic scale — i.e., there are roughly equal numbers of slow and fast fluctuators — the distribution of rates is $P(\gamma) \propto 1/\gamma$, and the total noise becomes:

\begin{equation}
    S(\omega) = \int_0^\infty S_k(\omega) P(\gamma)\, d\gamma \propto \int_0^\infty \frac{1}{\omega^2 + \gamma^2}\, d\gamma = \frac{\pi}{2\omega}
\end{equation}

This integral evaluates to $\pi/2\omega$, giving $S(\omega) \propto 1/\omega$. The $1/f$ spectrum is therefore not a property of any single fluctuator — it is an emergent property of having many fluctuators with a broad distribution of switching rates \cite{Paladino2014}.

\subsection{Mathematical Model}

The generalized $1/f^\alpha$ noise model is:

\begin{equation}
    \boxed{S_{1/f}(\omega) = \frac{A}{\omega^\alpha}}
\end{equation}

where $A$ is the noise amplitude and $\alpha$ is the spectral exponent. For pure $1/f$ noise $\alpha = 1$. In superconducting qubits, $\alpha$ typically lies in the range $0.7 \lesssim \alpha \lesssim 1.3$ depending on the device and dominant noise source \cite{Yoshihara2006, Bylander2011}.

This model diverges at $\omega = 0$, reflecting the fact that there is always a low-frequency cutoff set by the finite duration of the experiment. In numerical implementation, the value at $\omega = 0$ is set to zero.

\subsection{Implications for Filter Functions}

The Ramsey filter function $F_\text{Ramsey}(\omega, t) = 4\sin^2(\omega t/2)/\omega^2$ peaks at $\omega = 0$ and scales as $t^2$ there. Combined with $S_{1/f}(\omega) \propto 1/\omega$, the integrand $S(\omega)F(\omega,t)$ is heavily weighted toward low frequencies. This is why Ramsey dephasing in superconducting qubits is dominated by $1/f$ noise and why Ramsey $T_2^*$ times are typically much shorter than echo $T_2$ times \cite{Bylander2011}.

\section{White Noise Floor}

\subsection{Physical Origin}

At frequencies above the $1/f$ knee, the fluctuator-dominated noise falls below a constant background arising from Johnson-Nyquist thermal noise. Any resistive element at finite temperature $T$ generates voltage fluctuations with a flat power spectral density:

\begin{equation}
    S_V(\omega) = 4k_B T R
\end{equation}

independent of frequency \cite{Clerk2010}. In a superconducting qubit experiment, this noise enters through control lines, amplifiers, bias resistors, and readout electronics — all of which are at temperatures above the qubit's millikelvin operating point. The result is a frequency-independent noise floor that sets the minimum achievable $S(\omega)$.

\subsection{Mathematical Model}

\begin{equation}
    \boxed{S_\text{white}(\omega) = S_0}
\end{equation}

where $S_0$ is a constant. The crossover frequency $\omega_c$ at which $1/f$ noise and white noise contribute equally is found by setting $A/\omega_c^\alpha = S_0$, giving $\omega_c = (A/S_0)^{1/\alpha}$.

\subsection{Implications for Filter Functions}

Echo sequences suppress low-frequency $1/f$ noise by refocusing slow phase accumulation. However, white noise cannot be refocused — it contributes uniformly regardless of the pulse sequence. This sets a fundamental limit on coherence time achievable through dynamical decoupling: even with infinitely many pulses, the white noise floor remains and causes irreversible dephasing \cite{Cywinski2008}.

\section{Lorentzian TLS Peaks}

\subsection{Physical Origin}

Individual two-level system defects — atomic tunneling systems in amorphous oxides, dangling bonds at interfaces, or trapped charges — each have a well-defined resonance frequency $\omega_k$ and a relaxation rate $\gamma_k$ \cite{Muller2019, Martinis2005}. When a TLS couples to the qubit, it generates fluctuations at its resonance frequency, producing a peaked noise contribution. The lineshape is Lorentzian because the TLS is a damped two-state system:

\begin{equation}
    S_{TLS,k}(\omega) = \frac{A_k \gamma_k^2}{(\omega - \omega_k)^2 + \gamma_k^2}
\end{equation}

A slowly relaxing TLS ($\gamma_k$ small) produces a sharp, narrow peak. A fast-relaxing TLS produces a broad, low peak. In real devices, multiple TLS defects contribute at different frequencies, producing a forest of Lorentzian peaks superimposed on the $1/f$ + white background \cite{Muller2019}.

\subsection{Mathematical Model}

\begin{equation}
    \boxed{S_{TLS}(\omega) = \frac{A \gamma^2}{(\omega - \omega_0)^2 + \gamma^2}}
\end{equation}

where $A$ is the peak amplitude, $\omega_0$ is the TLS resonance frequency, and $\gamma$ is the half-width at half-maximum (HWHM), equal to the TLS relaxation rate.

\subsection{Implications for Filter Functions and Resonance Condition}

When a dynamical decoupling sequence's filter peak coincides with a TLS resonance — i.e., when $\omega_\text{peak}(n, t) = \omega_k$ — the coherence decay is dramatically enhanced rather than suppressed. This is the \textbf{resonance condition}:

\begin{equation}
    \frac{n\pi}{t} = \omega_k
\end{equation}

For CPMG-$n$ with varying $t$, the filter peak sweeps through the spectrum and periodically resonates with TLS defects, producing characteristic dips in the coherence as a function of $t$. This effect has been experimentally observed in superconducting flux qubits \cite{Bylander2011} and is one of the primary signatures used to identify and locate TLS defects in device characterization.

Crucially, this same resonance condition is what makes high-$n$ CPMG sequences useful for TLS spectroscopy: by sweeping $n$ and $t$, one can map out the locations of TLS resonances in $S(\omega)$ with high frequency resolution.

\section{Combined Noise Model}

The full noise model used in this project is:

\begin{equation}
    S(\omega) = \frac{A_{1/f}}{\omega^\alpha} + S_0 + \sum_k \frac{A_k \gamma_k^2}{(\omega - \omega_k)^2 + \gamma_k^2}
\end{equation}

This model is implemented in \texttt{simulation/noise\_models.py} and used as the ground-truth $S(\omega)$ for simulation validation of the inversion pipeline. The inversion problem — recovering $S(\omega)$ from measured $\chi(t)$ — is treated in documents 05--08.

\begin{thebibliography}{9}

\bibitem{Paladino2014}
E. Paladino, Y. M. Galperin, G. Falci, and B. L. Altshuler,
\textit{$1/f$ noise: Implications for solid-state quantum information},
Rev. Mod. Phys. \textbf{86}, 361 (2014).

\bibitem{Yoshihara2006}
F. Yoshihara, K. Harrabi, A. O. Niskanen, Y. Nakamura, and J. S. Tsai,
\textit{Decoherence of Flux Qubits due to $1/f$ Flux Noise},
Phys. Rev. Lett. \textbf{97}, 167001 (2006).

\bibitem{Bylander2011}
J. Bylander et al.,
\textit{Noise spectroscopy through dynamical decoupling with a superconducting flux qubit},
Nature Physics \textbf{7}, 565 (2011).

\bibitem{Clerk2010}
A. A. Clerk, M. H. Devoret, S. M. Girvin, F. Marquardt, and R. J. Schoelkopf,
\textit{Introduction to quantum noise, measurement, and amplification},
Rev. Mod. Phys. \textbf{82}, 1155 (2010).

\bibitem{Muller2019}
C. M\"{u}ller, J. H. Cole, and J. Lisenfeld,
\textit{Towards understanding two-level-systems in amorphous solids},
Rep. Prog. Phys. \textbf{82}, 124501 (2019).

\bibitem{Martinis2005}
J. M. Martinis et al.,
\textit{Decoherence in Josephson Qubits from Dielectric Loss},
Phys. Rev. Lett. \textbf{95}, 210503 (2005).

\bibitem{Cywinski2008}
L. Cywi\'{n}ski, R. M. Lutchyn, C. P. Nave, and S. Das Sarma,
\textit{How to enhance dephasing time in superconducting qubits},
Phys. Rev. B \textbf{77}, 174509 (2008).

\end{thebibliography}

\end{document}